\documentclass[12pt,a4paper]{article}
\usepackage[utf8]{inputenc}
\usepackage[T5]{fontenc}
\usepackage[vietnamese]{babel}
\usepackage{amsmath, amssymb, amsfonts}
\usepackage{enumitem}

\newcommand{\R}{\mathbb{R}}
\newcommand{\N}{\mathbb{N}}
\newcommand{\ip}[2]{\left\langle #1,#2\right\rangle}

\begin{document}

\section*{NỘI DUNG ÔN TẬP CUỐI KỲ -- THỜI GIAN LÀM BÀI 90 PHÚT}
\textbf{SV chỉ được sử dụng tài liệu GIẤY}

\underline{Câu 1.} Trên $V=\R^5$ (hoặc $V=\R^6$), cho tập hợp
\[
W=\{\cdots\mid\cdots\}.
\]
Chứng minh rằng $W$ là không gian véc tơ con của $V$. Tìm hệ sinh, cơ sở và xác định số chiều cho $W$.

\underline{Câu 2.} Trên $V=\R^3$, cho tập hợp
\[
a=\{u_1=\cdots;u_2=\cdots;u_3=\cdots\}.
\]
\begin{enumerate}[label=\alph*/]
\item Chứng minh rằng $a$ là một cơ sở của $\R^3$.
\item Trực chuẩn hóa cơ sở $a$ thành cơ sở trực chuẩn $b$.
\item Tìm tọa độ của một véc tơ $v=\cdots\in\R^3$ theo cơ sở $a$.
\item Tìm tọa độ của một véc tơ $v=\cdots\in\R^3$ theo cơ sở $b$.
\end{enumerate}

\underline{Câu 3:} Trên $V=\R^4$, với tích vô hướng $\ip{\cdot}{\cdot}=\cdots$, cho hệ
$S=\{u_1=\cdots;u_2=\cdots;u_3=\cdots\}$. Dùng pp Gram--Schmidt để trực chuẩn hóa $S$.
Cho véc tơ $\lambda=\cdots$. Tìm tọa độ của véc tơ $\lambda$ theo cơ sở vừa trực chuẩn hóa được.

\underline{Câu 4.} Cho ma trận vuông thực, cấp $2$, $A=\cdots$. Hãy chéo hóa ma trận $A$, rồi tìm $A^m$, $\forall m\in\N$ và suy ra $A^{2026}$.

\begin{center}
----------HẾT----------
\end{center}

\end{document}
